\chapter{Quant Developer}\label{ch:in:quant-developer}

In fiscal 2025 the sourced finance employers filed 169 labour condition applications for jobs whose titles named a
quantitative developer or engineer. They filed 122 of them under the occupation code of financial quantitative analysts,
25 under software developers, and 22 under at least three other codes, none of which reached ten. The code is the
employer's choice, and it follows the job more than the title does: behind ``quant developer'' stand three different
jobs, the researcher's engineer, the desk's \omterm{def:in:quant-developer:library}{library developer} and the \omterm{def:in:quant-developer:trading}{trading-system developer}. This chapter separates
them, shows what each ships and how it is judged, and measures what the choice of code says about pay.

\begin{center}\footnotesize
\begin{tabular}{@{}p{2.4cm}p{3.3cm}p{3.3cm}p{3.3cm}@{}}
\toprule
\multicolumn{4}{@{}l}{\textbf{Role cards: the three quant developers}}\\
\midrule
 & \omterm{def:in:quant-developer:research}{research engineer} & \omterm{def:in:quant-developer:library}{library developer} & \omterm{def:in:quant-developer:trading}{trading-system developer}\\
\midrule
works with & researchers & \omterm{def:in:bank-quant:strat}{desk strategists} and \omterm{def:in:bank-quant:library}{library quants} & traders and engineers\\
ships & research infrastructure and production versions of signals & the pricing library's code & the trading system\\
horizon & days to weeks & releases & microseconds to days\\
reports to & head of research & head of quantitative analytics & head of trading technology\\
codes in filings & 13-2099.01, 15-1252 & 13-2099.01, 15-1252 & 15-1252\\
taught in & Book~7, ch.~1; Book~15, ch.~1 & Book~5, ch.~28; Book~6, ch.~29 & Book~13, ch.~1; Book~16, ch.~21\\
\bottomrule
\end{tabular}
\end{center}

Book~16, chapter~21, defines the quant developer as an engineer who turns a researcher's prototype into production code,
owns its performance and correctness, and knows enough of the model to change it safely, and places the role in embedded
or central teams. This chapter takes that definition as given and asks which of three products the engineer is building.

\section{The researcher's engineer}

\begin{definition}[Research engineer]\label{def:in:quant-developer:research}
A \emph{research engineer}\index{research engineer} is a quant developer who works inside a research group and builds
what researchers work with and what their results become: data pipelines, backtesters and simulation tools, the
computing that runs experiments, and the production versions of signals the group has decided to trade.
\end{definition}

The \omterm{def:in:quant-developer:research}{research engineer}'s customers are researchers, and the product is their speed. Book~7 (chapter~1) describes the
research process the engineer serves, and Book~15 the platforms (research, data and risk) that a firm builds for it:
the point-in-time data store, the backtester that reproduces a result a year later, the experiment tracker, the job
scheduler on the cluster. At a systematic fund or a platform, where many researchers share one infrastructure, the
\omterm{def:in:quant-developer:research}{research engineer}'s work is multiplied by every researcher who uses it; at a market maker, where research sits close to
trading, the engineer also carries a signal from the research code into the trading system and checks that the two agree.

The job is judged on two things. The first is correctness under reproduction: a backtest that cannot be rerun, or a
production signal that differs from its research version, destroys the evidence research exists to produce. The second
is time: the hours between an idea and its first honest test. A researcher's \omterm{def:in:quant-researcher:feedback}{feedback speed} (chapter~17) is limited by
the market; the \omterm{def:in:quant-developer:research}{research engineer} sets the other limit.

\section{The desk's library developer}

\begin{definition}[Library developer]\label{def:in:quant-developer:library}
A \emph{library developer}\index{library developer} is a quant developer who writes and maintains the code of a
pricing and risk library for the quants who design its models: the interfaces, data structures, performance and
numerical plumbing that turn a \omterm{def:in:bank-quant:library}{library quant}'s methods into a product every desk and risk system can call.
\end{definition}

The \omterm{def:in:quant-developer:library}{library developer} works beside the \omterm{def:in:bank-quant:library}{library quant} of chapter~18, and the line between them is a matter of
emphasis: the quant owns the mathematics of a model, the developer the code that makes it fast, stable and callable from
elsewhere. Book~5, chapter~28, builds a pricing library with a core in C++20 and Rust behind a Python interface, and the
split of work there is the split of the two jobs: the numerical method on one side; memory layout, threading, the
interface, the build and the regression suite on the other. Where one library serves every desk of a bank, the job sits
with the group that owns it; where pricing is part of the trading system, as for an options market maker (Book~11), it
merges into the \omterm{def:in:quant-developer:trading}{trading-system developer}'s.

The \omterm{def:in:quant-developer:library}{library developer} is judged like the library itself: no desk's numbers change unless someone meant them to, every
release reproduces the last one's results to a stated tolerance, and the library is fast enough for the risk runs that
call it millions of times a night.

\section{The trading-system developer}

\begin{definition}[Trading-system developer]\label{def:in:quant-developer:trading}
A \emph{trading-system developer}\index{trading-system developer} is a quant developer who builds and runs the systems
that trade: market-data handlers, the strategy engine, order entry, risk checks and the monitoring around them, measured
in microseconds and in what they do when something fails.
\end{definition}

The \omterm{def:in:quant-developer:trading}{trading-system developer} works closest to money. The tick-to-trade latency and the latency budget of Book~13,
chapter~1, are this developer's units; the exchange protocols of Book~10 and the network and hardware of Book~14 are the
developer's environment. The quantitative part of the job is not the model but its implementation under constraints: a
signal that must be computed in two microseconds, an inventory skew that must be applied on every quote, a risk check
that must reject an order before it leaves. At market makers and high-frequency firms these developers are often filed as
software engineers, and chapter~20 covers the engineering that surrounds them.

\begin{dated}{2025-09}{Quant developers in the filings and the survey}\label{dat:in:quant-developer:codes}
Fiscal 2025 labour condition applications at the sourced finance employers: 169 in the quant-developer title family (122
coded 13-2099.01, 25 coded 15-1252, 22 under codes with fewer than ten each); the systematic funds' cell, 12 applications
from 3 employers, median offered base \$235\,000 (10th--90th percentile \$195\,500--272\,500); every other employer
kind's cell comes from one or two employers and is suppressed. Across all titles: software developers (15-1252) 5\,531
applications, median \$155\,000; financial quantitative analysts (13-2099.01) 1\,294, median \$150\,000. Occupational
survey, May 2025, software developers: securities industry median \$163\,290 (10th--90th \$104\,600--245\,380; 34\,940
employed); banks \$136\,620; software publishers \$164\,550 (173\,080 employed).
\end{dated}

\section{What each ships and how it is judged}

The three jobs share a title and much of a toolkit (a systems language, Python, version control, testing, some
statistics), but not a product or a clock.
\begin{center}\footnotesize
\begin{tabular}{@{}p{2.7cm}p{3.3cm}p{3.1cm}p{3.3cm}@{}}
\toprule
 & \omterm{def:in:quant-developer:research}{research engineer} & \omterm{def:in:quant-developer:library}{library developer} & \omterm{def:in:quant-developer:trading}{trading-system developer}\\
\midrule
ships & pipelines, tools, production signals & library releases & system releases, configurations\\
judged on & reproducibility; researchers' time & stability; speed; coverage & latency; uptime; correct behaviour\\
failure looks like & a result no one can rerun & all desks' numbers move & a loss in seconds\\
on call & rarely & at release & while markets are open\\
\bottomrule
\end{tabular}
\end{center}

The occupation codes follow the same line. The financial quantitative analysts' definition, ``Develop mathematical or
statistical models for risk management, asset optimization, pricing, or relative value analysis'', fits the \omterm{def:in:quant-developer:research}{research
engineer} and the \omterm{def:in:quant-developer:library}{library developer}; the software developers', ``Research, design, and develop computer and network
software'', fits the \omterm{def:in:quant-developer:trading}{trading-system developer}. O*NET places the first in its highest preparation band (most need graduate
school) and the second in the band below (most need a bachelor's degree). The code also moves: in 2018 the classification
merged the two software-developer codes of 2010 (15-1132 and 15-1133) into 15-1252, so fiscal 2021 filings need a
crosswalk before they can be compared with fiscal 2025's.

\begin{omfigure}
\begin{tikzpicture}
\begin{axis}[xbar stacked,width=10.5cm,height=5.2cm,xmin=0,xmax=100,ymin=0.4,ymax=4.6,bar width=11pt,ytick=data,
  yticklabels from table={figdata/industry/19-quant-developer/soc_mix.csv}{family},yticklabel style={font=\scriptsize},
  xticklabel style={font=\footnotesize},xlabel={\small share of fiscal 2025 filings, \%},table/col sep=comma,
  legend style={legend cell align=left,font=\scriptsize,at={(0.5,-0.32)},anchor=north,/tikz/every even column/.append style={column sep=6pt}},
  legend columns=5]
\addplot[fill=omThm,draw=none] table[x=quant,y=pos]{figdata/industry/19-quant-developer/soc_mix.csv};
\addplot[fill=omDef,draw=none] table[x=software,y=pos]{figdata/industry/19-quant-developer/soc_mix.csv};
\addplot[fill=omDef!40,draw=none] table[x=stats,y=pos]{figdata/industry/19-quant-developer/soc_mix.csv};
\addplot[fill=omThm!40,draw=none] table[x=data,y=pos]{figdata/industry/19-quant-developer/soc_mix.csv};
\addplot[fill=black!25,draw=none] table[x=other,y=pos]{figdata/industry/19-quant-developer/soc_mix.csv};
\legend{13-2099.01,15-1252,15-2041 or 15-2031,15-2051,other}
\end{axis}
\end{tikzpicture}

\omcaption{Occupation codes within four title families, fiscal 2025 filings at the sourced finance employers. Codes with
fewer than ten filings in a family are counted as ``other''. The quant-developer titles are filed mostly as quantitative
analysts, the software-engineer titles almost only as software developers. Data: \texttt{data/industry/lca\_titles.csv}
and \texttt{lca\_ranges.csv}, through \texttt{in\_quantdev.soc\_mix}.}\label{fig:in:quant-developer:mix}
\end{omfigure}

Does the code matter for pay? Across all titles at the sourced employers in fiscal 2025, filings coded as software
developers have a median offered base \$5\,000 above those coded as quantitative analysts, with a bootstrap interval from
$-\$372$ to \$5\,688 that just includes zero. Within each \omterm{def:in:pay-levels-by-role-firm-type-and-seniority:soc}{wage level} the sign reverses: the analysts' median is
higher at levels III and IV, by \$11\,000 and \$27\,092 (\cref{fig:in:quant-developer:levels,fig:in:quant-developer:gap}).
The two statements are consistent because the codes are used at different levels: half the software-developer filings
are at level IV, against one in seven of the analysts' (\cref{fig:in:quant-developer:levelmix}).

\begin{omfigure}
\begin{tikzpicture}
\begin{axis}[width=10.5cm,height=5.2cm,xmin=0.7,xmax=4.3,ymin=80,ymax=215,xtick={1,2,3,4},xticklabels={I,II,III,IV},
  xlabel={\small wage level},ylabel={\small median offered base, \$ thousand},tick label style={font=\footnotesize},
  grid=major,grid style={gray!20},table/col sep=comma,
  legend style={font=\scriptsize,at={(0.03,0.97)},anchor=north west,legend cell align=left}]
\addplot[omDef,thick,mark=*,error bars/.cd,y dir=both,y explicit]
  table[x expr=\thisrow{level}-0.05,y=dev,y error plus expr=\thisrow{dev_hi}-\thisrow{dev},
  y error minus expr=\thisrow{dev}-\thisrow{dev_lo}]{figdata/industry/19-quant-developer/levels.csv};
\addplot[omThm,thick,mark=square*,error bars/.cd,y dir=both,y explicit]
  table[x expr=\thisrow{level}+0.05,y=quant,y error plus expr=\thisrow{quant_hi}-\thisrow{quant},
  y error minus expr=\thisrow{quant}-\thisrow{quant_lo}]{figdata/industry/19-quant-developer/levels.csv};
\legend{software developers (15-1252),quantitative analysts (13-2099.01)}
\end{axis}
\end{tikzpicture}

\omcaption{Median offered base by \omterm{def:in:pay-levels-by-role-firm-type-and-seniority:soc}{wage level} for the two codes a quant developer is filed under, all titles at the sourced
finance employers, fiscal 2025, with 95\% bootstrap intervals for the median. Data: \texttt{data/industry/lca\_soc.csv},
through \texttt{in\_quantdev.soc\_cells}.}\label{fig:in:quant-developer:levels}
\end{omfigure}

\begin{omfigure}
\begin{tikzpicture}
\begin{axis}[width=10.5cm,height=5cm,xmin=0.5,xmax=5.5,ymin=-40,ymax=15,xtick={1,2,3,4,5},
  xticklabels={all levels,I,II,III,IV},xlabel={\small wage level},ylabel={\small gap, \$ thousand},
  tick label style={font=\footnotesize},grid=major,grid style={gray!20},table/col sep=comma,
  legend style={font=\scriptsize,at={(0.97,0.97)},anchor=north east,legend cell align=left}]
\addplot[only marks,mark=o,black!60,thick,error bars/.cd,y dir=both,y explicit]
  table[x expr=\thisrow{pos}-0.1,y=d2021,y error plus expr=\thisrow{hi2021}-\thisrow{d2021},
  y error minus expr=\thisrow{d2021}-\thisrow{lo2021}]{figdata/industry/19-quant-developer/gap.csv};
\addplot[only marks,mark=*,omDef,thick,error bars/.cd,y dir=both,y explicit]
  table[x expr=\thisrow{pos}+0.1,y=d2025,y error plus expr=\thisrow{hi2025}-\thisrow{d2025},
  y error minus expr=\thisrow{d2025}-\thisrow{lo2025}]{figdata/industry/19-quant-developer/gap.csv};
\draw[black!50] (axis cs:0.5,0) -- (axis cs:5.5,0);
\legend{fiscal 2021,fiscal 2025}
\end{axis}
\end{tikzpicture}

\omcaption{Median offered base of software-developer filings minus that of quantitative-analyst filings, overall and by
\omterm{def:in:pay-levels-by-role-firm-type-and-seniority:soc}{wage level}, with 95\% bootstrap intervals (fiscal 2021 codes merged by the 2018 crosswalk). At or above zero overall, negative or
indistinguishable from zero within each level. Data: \texttt{data/industry/lca\_soc\_gap.csv}, through
\texttt{in\_quantdev.gaps}.}\label{fig:in:quant-developer:gap}
\end{omfigure}

\begin{omfigure}
\begin{tikzpicture}
\begin{axis}[ybar,area legend,width=10.5cm,height=4.8cm,xmin=0.5,xmax=4.5,ymin=0,ymax=60,bar width=14pt,xtick={1,2,3,4},
  xticklabels={I,II,III,IV},xlabel={\small wage level},ylabel={\small share of filings, \%},
  tick label style={font=\footnotesize},grid=major,grid style={gray!20},table/col sep=comma,
  legend style={legend cell align=left,font=\scriptsize,at={(0.5,-0.42)},anchor=north,/tikz/every even column/.append style={column sep=8pt}},
  legend columns=2]
\addplot[fill=omDef,draw=none] table[x=level,y=dev]{figdata/industry/19-quant-developer/level_mix.csv};
\addplot[fill=omThm,draw=none] table[x=level,y=quant]{figdata/industry/19-quant-developer/level_mix.csv};
\legend{software developers (15-1252),quantitative analysts (13-2099.01)}
\end{axis}
\end{tikzpicture}

\omcaption{How the two codes' fiscal 2025 filings spread over the four \omterm{def:in:pay-levels-by-role-firm-type-and-seniority:soc}{wage levels} (filings with a stated level). Half the
software-developer filings are at level IV, which lifts their overall median. Data: \texttt{in\_quantdev.level\_mix}.}\label{fig:in:quant-developer:levelmix}
\end{omfigure}

The comparison says what the codes do not: that a developer is paid more for being called a software developer. At the
same stated level, the analysts' code carries the higher offer; overall, the software code carries more senior filings,
and most of them come from the banks' engineering workforces (4\,790 of the 5\,531, from eight banks). A quant developer's pay depends on the employer and
the level, as chapter~14 showed; the code is a description, and the two gaps are a warning about comparing medians over
populations with different mixes.

\section{Moving between the three}

The three jobs are close enough for people to move between them, and far enough for the move to cost something.
\begin{itemize}
\item \textbf{Trading-system to research engineering} brings a developer nearer to research: the systems skills carry
over, and a developer who has run the production version of a signal knows where research code breaks.
\item \textbf{Research engineering to research} is a change of product, from tools to results. A \omterm{def:in:quant-developer:research}{research engineer} who
already writes signals' production code can take a research problem of their own; firms differ in whether they allow it
(chapter~28).
\item \textbf{Library development to the desk} follows the bank's path of chapter~18: a \omterm{def:in:quant-developer:library}{library developer} who knows the
models becomes a \omterm{def:in:bank-quant:strat}{desk strategist}; one who prefers systems moves to trading technology.
\item \textbf{Between employers} the code and the title travel badly: a ``quant developer'' at a bank and at a market
maker may do the library and the trading-system job respectively. A candidate should ask which of the three products the
job ships and what the on-call duty is.
\end{itemize}

\section{Tutorial: one title, several codes}\label{tut:in:quant-developer:tut}

\begin{tutorial}
\textbf{Goal.} Separate the title from the job and from the code, and measure what the code says about pay.
\textbf{End state:} \cref{fig:in:quant-developer:mix,fig:in:quant-developer:levels,fig:in:quant-developer:gap,fig:in:quant-developer:levelmix}.
\begin{tutsteps}
\item \textbf{Normalise titles.} \texttt{firm.roles.normalise\_title} strips teams and locations, takes out seniority
words, expands abbreviations and drops grade numerals (\cref{lst:in:quant-developer:normalise}); on the chapter's
illustrative list, ``Sr.\ Quant Developer - Equities'', ``VP, Quant Dev'' and ``Quant Dev L3'' all become
\texttt{QUANTITATIVE DEVELOPER}.
\omcode{code/firm/roles/firm_roles.py}{224}{243}{The job-title normaliser.}\label{lst:in:quant-developer:normalise}
\item \textbf{Cross-tabulate.} \texttt{crosstab(pairs)} counts title families by occupation code and suppresses cells under
ten; chapter~14's \texttt{lca\_titles.csv} is that table for the filings.
\item \textbf{Cells by code and level.} \texttt{in\_lca\_soc\_derive.py} rebuilds the code tables from chapter~14's cached
records (no second pass over the raw files), with the three-employer rule.
\item \textbf{The gap.} \texttt{median\_gap(x, y, rng)} gives the difference and ratio of medians with independent
bootstrap intervals (\cref{lst:in:quant-developer:gap}).
\omcode{code/firm/roles/firm_roles.py}{254}{265}{The gap between two medians with bootstrap intervals.}\label{lst:in:quant-developer:gap}
\end{tutsteps}
\end{tutorial}

For fiscal 2025: overall gap \$5\,000 ($-\$372$ to \$5\,688); level III $-\$11\,000$ ($-\$19\,038$ to $-\$6\,600$); level
IV $-\$27\,092$ ($-\$37\,500$ to $-\$19\,900$). Holding the level mix at the analysts', the software developers' median by
level averages \$7\,789 below the analysts'.

\textbf{What to change next.} Take the employer into account (a gap within each employer, then averaged); separate
full-time new hires from transfers; match the normaliser's cleaned titles to the codes directly, which needs the titles
that this chapter's single pass over the files did not keep.

\section{Build: title normaliser and code comparison}\label{bld:in:quant-developer:roles}

\begin{build}
\bfield{Purpose} Add the three quant developers to the role registry, and the tools to study titles and codes.
\bfield{Interface} \texttt{firm.roles}: the cards \texttt{research engineer}, \texttt{library developer},
\texttt{trading-system developer}; \texttt{normalise\_title(title) -> (clean, seniority)};
\texttt{crosstab(pairs, min\_cell)}; \texttt{median\_gap(x, y, rng, n\_boot, min\_cell)}.
\bfield{Rules} Seniority is recorded, not kept in the title; the first part of a title that names a job wins; cells
under ten are suppressed; a gap is not computed from fewer than ten values on either side.
\bfield{Acceptance tests} \texttt{code/firm/roles/tests/}: six titles normalised to their expected forms; a cross-tab
suppresses a cell of three; a planted gap of 50 lies inside its interval; a population of nine gives no gap.
\bfield{Stretch} Learn the normaliser's abbreviation table from data; stratify the gap by employer; a permutation test
of the gap.
\end{build}

\begin{omsources}
\item O*NET OnLine, occupations 15-1252, 13-2099.01, 15-2041 and 15-2051 (CC BY 4.0).
\item Bureau of Labor Statistics, SOC 2010 to 2018 crosswalk; occupational survey, May 2025.
\item Chapter~14's tables from the Department of Labor's LCA files.
\item Book~16, chapter~21, for the quant developer and embedded and central teams.
\end{omsources}

\section{Exercises}

\begin{exercise}[$\star$]\label{exo:in:quant-developer:1}
What share of the fiscal 2025 quant-developer filings were coded as financial quantitative analysts, and what share as
software developers?
\end{exercise}

\begin{exercise}[$\star$]\label{exo:in:quant-developer:2}
What does \texttt{normalise\_title} return for ``Sr.\ Quant Dev (Rates), NY''?
\end{exercise}

\begin{exercise}[$\star$]\label{exo:in:quant-developer:3}
From the dated box, how much higher is the survey's median for software developers in the securities industry than in
banks, as a ratio?
\end{exercise}

\begin{exercise}[$\star\star$]\label{exo:in:quant-developer:4}
Explain in three sentences how the software developers' median can exceed the analysts' overall while falling below it at
levels III and IV.
\end{exercise}

\begin{exercise}[$\star\star$]\label{exo:in:quant-developer:5}
Weight each code's median by level with the analysts' level shares (7.8\%, 51.1\%, 26.7\%, 14.4\%). What is the
difference?
\end{exercise}

\begin{exercise}[$\star\star$]\label{exo:in:quant-developer:6}
A graduate wants to become a researcher within three years. Which of the three quant-developer jobs gives the best
start, and what should she ask at interview?
\end{exercise}

\begin{exercise}[$\star\star\star$]\label{exo:in:quant-developer:7}
\emph{Coding.} Write a function that calls a gap distinguishable from zero when its interval excludes zero, and apply it to
the ten gaps of \texttt{lca\_soc\_gap.csv}. How many are distinguishable?
\end{exercise}

\begin{exercise}[$\star\star\star$]\label{exo:in:quant-developer:8}
\emph{Find the flaw.} ``Software developers at finance employers earn more than quantitative analysts: their median is
\$5\,000 higher. A quant developer should ask to be filed as a software developer.''
\end{exercise}

\section{Problem: One Title, Three Jobs}

\begin{problem}[{Weekend problem --- one title, three jobs}]\label{pb:in:quant-developer:1}
A developer has three offers, each titled ``quantitative developer'': at a bank's analytics group, at a systematic fund's
research group and at a market maker's trading-technology group. He wants to know what the title hides.

\medskip\noindent\textbf{Part I --- The jobs.}
\begin{enumerate}
\item Define the quant developer (Book~16) and the three jobs.
\item What does each ship?
\item What is each judged on, and what does failure look like?
\item Which job carries on-call duty while markets are open?
\item Where does each sit in the organisation?
\end{enumerate}

\medskip\noindent\textbf{Part II --- The codes.}
\begin{enumerate}[resume]
\item What are occupation codes 15-1252 and 13-2099.01, and which jobs fit each?
\item What preparation does O*NET attach to each?
\item What changed in the software-developer codes in 2018?
\item State the SOC mix of the fiscal 2025 quant-developer filings.
\item Why are most employer kinds' pay cells for the family suppressed?
\end{enumerate}

\medskip\noindent\textbf{Part III --- The gap.}
\begin{enumerate}[resume]
\item State the method of \texttt{median\_gap}.
\item Give the overall fiscal 2025 gap and its interval.
\item Give the gaps at levels III and IV.
\item Explain the reversal.
\item What happens to the gap when the level mix is held at the analysts'?
\end{enumerate}

\medskip\noindent\textbf{Part IV --- The verdict.}
\begin{enumerate}[resume]
\item State the \emph{named result}: the SOC mix of quant-developer filings and the difference in median offered base
between the software-developer and quantitative-analyst codes, with its interval.
\item Does the code he is filed under matter for his pay?
\item What does the systematic funds' cell show, and why is it the only one?
\item Which question separates the three offers best?
\item In two sentences, answer him.
\end{enumerate}
\end{problem}

\section{Interview questions}

\begin{interviewq}[$\star$ \iqroles{developer}]\label{iq:in:quant-developer:1}
A researcher's Python signal and your C++ version disagree in the fourth decimal on 0.1\% of days. What do you do?
\end{interviewq}

\begin{interviewq}[$\star$ \iqroles{developer}]\label{iq:in:quant-developer:2}
How would you make a backtest reproducible a year later?
\end{interviewq}

\begin{interviewq}[$\star\star$ \iqroles{developer, bank}]\label{iq:in:quant-developer:3}
A pricing function is called ten million times in the nightly risk run and takes 40 microseconds. Where do you look first?
\end{interviewq}

\begin{interviewq}[$\star\star$ \iqroles{developer, trader}]\label{iq:in:quant-developer:4}
The trading system's median tick-to-trade is fine but the 99.9th percentile doubled after a release. How do you find
out why?
\end{interviewq}

\begin{interviewq}[$\star\star$ \iqroles{developer, researcher}]\label{iq:in:quant-developer:5}
What would you put in a signal's interface so that research and production cannot silently diverge?
\end{interviewq}

\begin{interviewq}[$\star\star\star$ \iqroles{developer}]\label{iq:in:quant-developer:6}
Design the regression test suite for a pricing library that many desks use.
\end{interviewq}
